\section*{Hoofdstuk \ref{ch:b1:crystals-ionic-covalent} --- Kristallen II: ionaire, covalente en moleculaire vaste stoffen}

\begin{solution}{exo:b1:crystals-ionic-covalent:1}
Anionen: $8 \times \frac18 + 6 \times \frac12 = 4$; kationen: $12 \times
\frac14 + 1 = 4$. Elk ion heeft 6 buren van de andere soort (6:6). Vier
kationen en vier anionen: \ce{NaCl}.
\end{solution}

\begin{solution}{exo:b1:crystals-ionic-covalent:2}
De halve lichaamsdiagonaal: $a\sqrt3/2 = 412.3 \times 0.866 = \qty{357.1}{pm}$.
Eén \ce{Cs+} en $8 \times \frac18 = 1$ \ce{Cl-}: één formule-eenheid.
\end{solution}

\begin{solution}{exo:b1:crystals-ionic-covalent:3}
$x = 102/181 = 0.56$, tussen 0.414 en 0.732: octaëdrische coördinatie
(6), het \omterm{def:b1:crystals-ionic-covalent:structure-type}{steenzouttype}.
\end{solution}

\begin{solution}{exo:b1:crystals-ionic-covalent:4}
Metallisch: koper. Ionair: natriumchloride, calciumfluoride. Covalent:
diamant, kwarts. Moleculair: dijood, ijs.
\end{solution}

\begin{solution}{exo:b1:crystals-ionic-covalent:5}
Langs een ribbe zijn anion, kation en anion in contact: $a = 2(r_+ + r_-)$.
Langs een zijvlakdiagonaal liggen twee anionen op de afstand
$a\sqrt2/2$; die mogen niet overlappen: $a\sqrt2/2 \ge 2r_-$, dus
$(r_+ + r_-)\sqrt2 \ge 2r_-$, oftewel
$r_+/r_- \ge \sqrt2 - 1$.
\end{solution}

\begin{solution}{exo:b1:crystals-ionic-covalent:6}
$\rho = 4 \times 58.44/(6.022 \times 10^{23} \times (5.641 \times
10^{-8})^3) = \qty{2.16}{g/cm^3}$, tegen \qty{2.17}{g/cm^3} gemeten.
\end{solution}

\begin{solution}{exo:b1:crystals-ionic-covalent:7}
$\rho = 168.36/(6.022 \times 10^{23} \times (4.123 \times 10^{-8})^3) =
\qty{3.99}{g/cm^3}$, de gemeten waarde.
\end{solution}

\begin{solution}{exo:b1:crystals-ionic-covalent:8}
\ce{Zn-S} $= a\sqrt3/4 = \qty{234.2}{pm}$;
$\rho = 4 \times 97.44/(6.022 \times 10^{23} \times (5.409 \times
10^{-8})^3) = \qty{4.09}{g/cm^3}$, dicht bij de gemeten
\qty{4.04}{g/cm^3} van natuurlijk sfaleriet.
\end{solution}

\begin{solution}{exo:b1:crystals-ionic-covalent:9}
\ce{C-C} $= 245.6/\sqrt3 = \qty{141.8}{pm}$; afstand tussen de lagen
$669.6/2 =
\qty{334.8}{pm}$. Celvolume $\frac{\sqrt3}{2}a^2c = \qty{3.498e-23}{cm^3}$,
$\rho = 4 \times 12.01/(6.022 \times 10^{23} \times 3.498 \times 10^{-23})
= \qty{2.28}{g/cm^3}$, tegen \qty{3.52}{g/cm^3} voor diamant: de lagen
liggen ver uit elkaar en worden alleen door londonkrachten
bijeengehouden.
\end{solution}

\begin{solution}{exo:b1:crystals-ionic-covalent:10}
Zie het bewijs van \cref{prop:b1:crystals-ionic-covalent:diamond}:
$C = \pi\sqrt3/16 = 0.34$. De hardheid komt van de sterkte en de
driedimensionale starheid van de \omterm{def:b1:lewis-resonance-vsepr:lewis}{covalente bindingen}, niet van de
stapeling: elk atoom heeft maar vier buren, onder vaste hoeken, wat veel
lege ruimte overlaat.
\end{solution}

\begin{solution}{exo:b1:crystals-ionic-covalent:11}
\ce{Si-Si} $= a\sqrt3/4 = \qty{235.2}{pm}$, \omterm{def:b1:periodicity:radius}{covalente straal}
\qty{117.6}{pm}; $\rho = 8 \times 28.09/(6.022 \times 10^{23} \times
(5.431 \times 10^{-8})^3) = \qty{2.33}{g/cm^3}$, de gemeten waarde.
\end{solution}

\begin{solution}{exo:b1:crystals-ionic-covalent:12}
$x = 152/181 = 0.84 > 0.732$: de regel voorspelt het \omterm{def:b1:crystals-ionic-covalent:structure-type}{cesiumchloridetype}.
In de steenzoutstructuur is \ce{Rb-Cl} $= a/2 = \qty{329.1}{pm}$, dicht
bij $152 + 181 = \qty{333}{pm}$: de waargenomen structuur is consistent.
De regel behandelt ionen als harde bollen en verwaarloost hun
\omterm{def:b1:periodicity:polarisability}{polariseerbaarheid} en het kleine energieverschil tussen de twee
structuren (rubidiumchloride neemt onder druk het \omterm{def:b1:crystals-ionic-covalent:structure-type}{cesiumchloridetype}
aan); hij is een leidraad, geen wet.
\end{solution}

\begin{solution}{pb:b1:crystals-ionic-covalent:1}
\textbf{1.} Op de 8 hoekpunten en de 6 zijvlakmiddelpunten: $8 \times \frac18 + 6
\times \frac12 = 4$ ionen.
\textbf{2.} 8 \omterm{def:b1:crystals-metals:site}{tetraëdrische holten}, alle bezet: 8 \ce{F-} voor 4
\ce{Ca^{2+}}, verhouding 1:2, \ce{CaF2}.
\textbf{3.} \ce{F-}: 4 \ce{Ca^{2+}} op de hoekpunten van een tetraëder.
\ce{Ca^{2+}}: 8 \ce{F-} op de hoekpunten van een kubus.
\textbf{4.} $a\sqrt3/4 = 546.3 \times 0.4330 = \qty{236.6}{pm}$.
\textbf{5.} $112 + 131 = \qty{243}{pm}$: binnen 3\,\%.
\textbf{6.} De \ce{F-}-ionen vormen een enkelvoudig kubische schikking met
ribbe $a/2 =
\qty{273.2}{pm}$, iets meer dan $2 \times 131 = \qty{262}{pm}$: ze raken
elkaar bijna.
\textbf{7.} $x = 112/131 = 0.85$.
\textbf{8.} $x > 0.732$: coördinatie 8.
\textbf{9.} Het \omterm{def:b1:crystals-ionic-covalent:structure-type}{cesiumchloridetype} heeft evenveel kationen als anionen;
met twee keer zoveel anionen kan maar de helft van de kubussen van
anionen een kation bevatten.
\textbf{10.} De 8 \ce{F-} van de cel vormen een kubus van 8 kleine
kubussen met ribbe $a/2$; de \ce{Ca^{2+}}-ionen bezetten afwisselend de
middelpunten van 4 ervan, zoals de zwarte velden van een
driedimensionaal schaakbord.
\textbf{11.} $C = \frac43\pi(4 \times 112^3 + 8 \times 131^3)/546.3^3 = 0.61$.
\textbf{12.} \ce{O^{2-}} op het \omterm{def:b1:crystals-metals:lattice}{FCC-rooster} (hoekpunten en
zijvlakmiddelpunten), \ce{Li+} in de 8 \omterm{def:b1:crystals-metals:site}{tetraëdrische holten}.
\textbf{13.} $a\sqrt3/4 = \qty{200.0}{pm}$; $59 + 142 = \qty{201}{pm}$.
\textbf{14.} $\rho = 4 \times 29.88/(6.022 \times 10^{23} \times (4.619
\times 10^{-8})^3) = \qty{2.01}{g/cm^3}$, de gemeten waarde.
\textbf{15.} Zinkblende: diamant is zinkblende met koolstof op beide
soorten posities.
\textbf{16.} $a\sqrt3/4 = \qty{154.5}{pm}$; $\rho = 8 \times
12.01/(6.022 \times 10^{23} \times (3.567 \times 10^{-8})^3) =
\qty{3.52}{g/cm^3}$.
\textbf{17.} In diamant is maar de helft van de \omterm{def:b1:crystals-metals:site}{tetraëdrische holten}
bezet, en dan nog door atomen die even groot zijn als die van het
\omterm{def:b1:crystals-metals:lattice}{rooster} en elkaar alleen langs de bindingen raken: $C = 0.34$. In
fluoriet zijn alle holten bezet en vullen de kleine kationen de ruimten
tussen grote anionen: $C = 0.61$.
\textbf{18.} $40.08 + 2 \times 19.00 = \qty{78.08}{g/mol}$.
\textbf{19.} $4 \times 78.08/(6.022 \times 10^{23}) = \qty{5.186e-22}{g}$.
\textbf{20.} $(5.463 \times 10^{-8})^3 = \qty{1.630e-22}{cm^3}$.
\textbf{21.} Eén ontbrekend \ce{F-} per honderd cellen neemt $19.00/(100
\times 312.3)$ van de massa weg, 0.06\,\%: onzichtbaar op drie cijfers (en
er moet ook een \ce{Ca^{2+}} verdwijnen, anders zou het \omterm{def:b1:crystals-metals:crystal}{kristal} niet
neutraal zijn).
\textbf{22.} $\rho = 5.186 \times 10^{-22}/1.630 \times 10^{-22} =
\textbf{\qty{3.18}{g/cm^3}}$, precies de gemeten dichtheid: het model van
de cel, opgebouwd uit één röntgenmeting van $a$, verklaart de massa van
het mineraal.
\end{solution}
